\documentclass[11pt]{article}

% review = anonymous submission build. Switch to "final" for camera-ready.
\usepackage[review]{acl}

\usepackage{times}
\usepackage{latexsym}
\usepackage[T1]{fontenc}
\usepackage[utf8]{inputenc}
\usepackage{microtype}
% \usepackage{inconsolata}  % optional typewriter font; enable on Overleaf
\usepackage{graphicx}
\usepackage{booktabs}

\title{TODO: Title (state the finding, not the method name)}

% Anonymous for review. Do not put names, affiliations, or a non-anonymized repo link here.
\author{Anonymous EACL submission}

\begin{document}
\maketitle

\begin{abstract}
TODO (write last). Three sentences: the problem, what you found, why it matters.
\end{abstract}

\section{Introduction}
% TODO (yours): the problem in your own words.
% Must contain: why outcome-only reward is a weak signal for multi-step reasoning;
% what TAPR does; the 3 to 4 concrete contributions, each backed by a table or figure below.

\section{Related Work}
% TODO (yours): process reward models, outcome vs process supervision, GRPO,
% LLM-as-judge for reasoning steps. Cite only papers you have opened.
% PROGRS and the neighbouring work from your notes go here.

\section{Method}
% TODO: TAPR reward definition, the judge prompts (isolated / transition / bare),
% lambda, and how the transition reward combines with the outcome reward.

\section{Experimental Setup}
% TODO: student (Qwen2.5-3B, 4-bit LoRA), judge, GRPO settings, tasks, seeds,
% evaluation sets and sizes (GSM8K full test n=1319; StrategyQA n=200), compute.

\section{Results}

\subsection{Judge validity against hand labels}
% Measured (14B judge, 40 hand-labelled traces over 20 problems, exploratory).
% TODO: add the 7B judge once scored. Intervals resample problems, not traces.
\begin{table}[t]
\centering
\small
\begin{tabular}{lcc}
\toprule
Scorer & Step AUC & 95\% CI \\
\midrule
bare       & 0.89 & [0.79, 0.97] \\
isolated   & 0.82 & [0.69, 0.94] \\
transition & 0.82 & [0.73, 0.90] \\
\bottomrule
\end{tabular}
\caption{Step-level AUC of the 14B judge against first-break hand labels
(19 breaks). TODO: 7B judge row.}
\label{tab:judge}
\end{table}

\subsection{Training: ORM vs TAPR}
% TODO: 3 seeds each, mean and std, full GSM8K (n=1319) plus StrategyQA transfer,
% paired significance on per-problem outputs.

\subsection{Reward degeneracy}
% TODO: within-group zero-variance rates for outcome vs TAPR reward,
% and whether the added variance tracks quality (variance-informativeness check).

\section{Analysis}
% TODO: policy drift (KL, LoRA norm, response length), and what changed in behaviour.

\section*{Limitations}
% REQUIRED by ACL/ARR. Be specific and honest. Candidates already known:
% 3B student only; two tasks; 40 hand labels, LLM-drafted and human-reviewed;
% 19 break steps; most breaks are plain arithmetic errors, so the population TAPR
% targets is small on GSM8K; judge accuracy measured on a small calibration set.

\section*{Ethics and Responsible NLP}
% TODO: complete the Responsible NLP checklist on submission.
% Disclose AI writing assistance as the ACL policy requires.

\bibliography{references}
\bibliographystyle{acl_natbib}

\appendix
\section{Hand-labelling protocol}
\label{sec:appendix}
% Rules are in notes/HANDLABEL_RULES.md. Summarize: first-break definition,
% four rules, break types, who labelled what (L01-L10 agreed, L11-L40 drafted and reviewed).

\end{document}
